% Part: propositional-logic
% Chapter: syntax-and-semantics
% Section: introduction

\documentclass[../../../include/open-logic-section]{subfiles}

\begin{document}

\olfileid{pl}{syn}{int}

\olsection{परिचय}

विधानीय तर्कशास्त्रात विधानीय संयोजके
$\lnot$, $\land$, $\lor$, $\lif$ आणि $\liff$ वापरून
!!{propositional variable}s पासून बनवलेल्या !!{formula}s चा विचार
केला जातो. अंतःप्रेरणेने पाहता, !!a{propositional variable}~$p$
हे सत्य किंवा असत्य असलेल्या एखाद्या वाक्याच्या किंवा विधानाच्या
जागी येते. !!a{formula} मधील !!{propositional variable} चे
``सत्यतामूल्य'' ठरले की, विधानीय संयोजके वापरून त्यांच्यापासून
बनवलेल्या कोणत्याही !!{formula}s चे सत्यतामूल्यही ठरते. विधानीय
तर्कशास्त्राला आपण \emph{सत्यता-फलनात्मक} म्हणतो, कारण त्याची
चिन्हार्थमीमांसा सत्यतामूल्यांच्या फलनांनी दिली जाते. विशेषतः,
विधानीय तर्कशास्त्रात सत्य आणि असत्य यांचे पुढील प्रकारचे निर्धारण
आपण विचारात घेत नाही: उदाहरणार्थ, एखादी गोष्ट केवळ परिस्थितीवश
सत्य असण्याऐवजी अनिवार्यपणे सत्य आहे का, ती सत्य आहे हे ज्ञात आहे
का, किंवा ती आत्ता सत्य आहे की पूर्वी सत्य होती अथवा पुढे सत्य
असेल. आपण फक्त सत्य ($\True$) आणि असत्य ($\False$) ही दोन
सत्यतामूल्ये विचारात घेतो. म्हणून एखादे विधान सत्यही नसते व
असत्यही नसते, किंवा ते केवळ अर्धसत्य असते, ही शक्यता चर्चेतून
वगळतो. तसेच, ज्यांच्यापासून बनवलेल्या !!a{formula} चे सत्यतामूल्य
त्याच्या भागांच्या सत्यतामूल्यांनी पूर्णपणे ठरते (त्याच्या अर्थाने,
उदाहरणार्थ, ठरत नाही), अशाच संयोजकांवर आपण लक्ष केंद्रित करतो.
विशेषतः, इंग्रजीतील सशर्त विधानांचे सत्यतामूल्य या अर्थाने
सत्यता-फलनात्मक असते की नाही, हा वादग्रस्त प्रश्न आहे. वास्तविक
अभिव्यंजन~$\lif$ सत्यता-फलनात्मक आहे; सत्यता-फलनात्मक नसलेल्या सशर्त
विधानांचा विचार इतर तर्कपद्धती करतात.

सत्यता-फलनात्मक विधानीय तर्कशास्त्राची उपपत्ती आणि अधिउपपत्ती
विकसित करण्यासाठी, आपण प्रथम त्याच्या व्यंजकांची विन्यासमीमांसा
आणि चिन्हार्थमीमांसा परिभाषित केली पाहिजे. संयोजके वापरून
!!{propositional variable}s पासून !!{formula}s तयार करण्याची एक
पद्धत आपण वर्णन करू. पर्यायी व्याख्या शक्य आहेत. इतर पद्धती वेगळी
चिन्हे निवडतील, संयोजकांचे वेगळे संच मूलभूत म्हणून निवडतील आणि
कंसांचा वेगळ्या रीतीने वापर करतील (किंवा तथाकथित पोलिश
संकेतपद्धतीप्रमाणे कंस अजिबात वापरणार नाहीत). तरी सर्व दृष्टिकोनांत
सामाईक बाब ही की रचनानियम !!{formula}s चा संच \emph{विगमनाने}
परिभाषित करतात. हे योग्य रीतीने केल्यास, रचनानियमांनुसार प्रत्येक
व्यंजक मूलतः एकाच रीतीने तयार होऊ शकते. त्यामुळे तयार होणारी
व्यंजके \emph{एकमेव रीतीने वाचनीय} असतात. परिणामी, त्यांची ही
विगमनाधारित व्याख्या आपल्याला त्याच पद्धतीने---विगमनाधारित
व्याख्येने---त्या व्यंजकांना अर्थ देता येतो, याची खात्री देते.

व्यंजकांना अर्थ देणे हा चिन्हार्थमीमांसेचा विषय आहे. विधानीय
तर्कशास्त्राच्या चिन्हार्थमीमांसेतील मध्यवर्ती संकल्पना म्हणजे
!!a{valuation} मधील पूर्ति. !!^a{valuation}~$\pAssign{v}$ हे
!!{propositional variable}s ना $\True$, $\False$ ही सत्यतामूल्ये
नेमून देते. कोणतेही !!{valuation} कोणत्याही !!{formula}~$!A$
साठी $\pValue{v}(!A)$ हे सत्यतामूल्य ठरवते. !!^a{formula} ची
!!a{valuation}~$\pAssign{v}$ मध्ये पूर्ति तेव्हा आणि केवळ तेव्हाच
होते जेव्हा $\pValue{v}(!A) = \True$; हे आपण $\pSat{v}{!A}$ असे
लिहितो. तार्किक संयोजकांची सत्यता-फलने वापरून, उदाहरणार्थ
$!A \land !B$ ची पूर्ति $!A$ आणि~$!B$ यांची पूर्ति होते (किंवा
होत नाही) यावरून परिभाषित करून, हा संबंधही~$!A$ च्या रचनेवरील
विगमनाने परिभाषित करता येतो.

वाक्यांसाठी असलेल्या $\pSat{v}{!A}$ या पूर्ति-संबंधाच्या आधारे
आपण सर्वतः सत्यता, तार्किक निष्पन्नता आणि पूर्ततायोग्यता या मूलभूत
चिन्हार्थविषयक संकल्पना परिभाषित करू शकतो. !!^a{formula} हे सर्वतः
सत्य सूत्र, $\Entails !A$, असते, जर प्रत्येक !!{valuation} त्याची
पूर्ति करत असेल, म्हणजे कोणत्याही~$\pAssign{v}$ साठी
$\pValue{v}(!A) = \True$ असेल. !!{formula}s च्या एखाद्या संचापासून,
$\Gamma \Entails !A$, ते तार्किकरीत्या निष्पन्न होते, जर~$\Gamma$
मधील सर्व !!{formula}s ची पूर्ति करणारे प्रत्येक !!{valuation}
~$!A$ चीही पूर्ति करत असेल. तसेच, !!{formula}s चा एखादा संच
पूर्ततायोग्य असतो, जर एखादे !!{valuation} त्यातील सर्व !!{formula}s
ची एकाच वेळी पूर्ति करत असेल. !!{formula}s विगमनाने
परिभाषित केलेली असल्यामुळे, आणि पूर्तिही !!{formula}s च्या रचनेवरील
विगमनाने परिभाषित केलेली असल्यामुळे, आपल्या चिन्हार्थमीमांसेचे
गुणधर्म सिद्ध करण्यासाठी आणि परिभाषित केलेल्या चिन्हार्थविषयक
संकल्पनांचा परस्परसंबंध दाखवण्यासाठी आपण विगमन वापरू शकतो.

\end{document}
